\setlength{\baselineskip}{20pt}
\chapter{引言}
\label{cha:intro}

\section{研究背景与意义}

随着计算机视觉技术的快速发展，图像判别式任务已成为人工智能领域中最活跃的研究方向之一。判别式任务涵盖图像分类\cite{agarap2017architecture,albawi2017understanding,deng2012mnist}、目标检测\cite{ren2015faster,he2017mask,redmon2018yolov3}、图像分割\cite{long2015fully,xie2021segformer}以及图像检索\cite{lin2015deep,weyand2020google}等广泛的应用场景，其核心目标在于从输入图像中学习有效的特征表示，进而完成对数据的精准预测与理解。在智能交通、公共安全、医疗影像和自动驾驶等实际应用中，判别式视觉模型的性能直接决定了系统的可靠性与可用性，因此受到了学术界和工业界的广泛关注。

表征学习是判别式视觉任务的核心基础，其被定义为``学习数据的表示（或特征），使得在构建分类器或其他预测器时能够更容易地提取有用信息''\cite{6472238}。一个典型的判别式视觉框架通常由负责特征提取的视觉骨干网络（如ResNet\cite{he2016deep}、ViT\cite{dosovitskiy2020image}）和面向特定任务的预测头（如用于分类的MLP\cite{krizhevsky2012imagenet}、用于检测的RCNN\cite{ren2015faster}、用于分割的FCN\cite{long2015fully}）组成。其中，骨干网络的特征表达能力对最终的任务性能起着决定性作用。

% 【建议插入图片：判别式视觉任务通用框架图】
% 内容：展示"输入图像 → 骨干网络(特征提取) → 特征表示 → 任务头 → 任务输出"的通用流水线，
% 并在特征表示处标注噪声干扰来源（输入噪声、传播噪声），形成对后文噪声瓶颈问题的直观铺垫。

近年来，深度学习技术的持续演进推动了判别式表征学习取得了显著进步。从早期的卷积神经网络（CNN）\cite{lecun1998gradient,krizhevsky2012imagenet}到残差网络ResNet\cite{he2016deep}，再到基于自注意力机制的Vision Transformer（ViT）\cite{dosovitskiy2020image,vaswani2017attention}和具有层级结构的Swin Transformer\cite{liu2021swin}，以及最新兴起的状态空间模型VMamba\cite{liu2024vmamba}，骨干网络的表达能力不断增强，在图像分类\cite{deng2009imagenet}、行人重识别\cite{ye2021deep,bedagkar2014survey}、目标检测\cite{he2017mask,carion2020end}和图像分割\cite{long2015fully,zhou2017scene}等任务上持续刷新性能记录。然而，尽管这些架构在模型容量和结构设计上取得了重大突破，特征中残留的噪声仍然是制约模型性能进一步提升的重要瓶颈。在推理阶段，噪声可能来源于多个方面：既包括输入图像本身携带的噪声（如光照变化、遮挡、背景杂乱等），也包括特征在网络逐层传播过程中累积的表示噪声。这些噪声导致类内特征分布发散、类间特征边界模糊，从而降低特征的判别性。

在此背景下，扩散模型（Diffusion Model）的兴起为上述问题提供了全新的思路。扩散模型\cite{sohl2015deep}的核心思想是通过向数据持续添加高斯噪声模拟前向扩散过程，再训练去噪网络学习反向去噪过程以恢复原始数据分布。自去噪扩散概率模型（DDPM）\cite{ho2020denoising}提出以来，扩散模型在图像生成\cite{rombach2022high,ramesh2021zero}、视频生成\cite{ho2022imagen}和分子生成\cite{hoogeboom2022equivariant}等领域均展现出卓越的生成能力，充分表明其具备优良的分布建模能力和强大的噪声抑制能力\cite{10419041}。这启发研究者思考：能否将扩散模型经过充分验证的去噪能力从生成式任务迁移至判别式任务，用于增强特征表示的质量？

综上所述，如何设计一种通用的特征增强范式，有效利用扩散模型的去噪能力来提升判别式表征的质量，同时避免推理阶段的额外计算开销，具有重要的理论意义和实际应用价值。一方面，从理论角度而言，探索生成式扩散模型与判别式表征学习的统一建模，有望为理解两类学习范式之间的内在联系提供新的视角，推动深度学习理论的发展。另一方面，从应用角度而言，设计零额外推理开销的特征去噪方法，能够在不改变已有模型部署架构的前提下增量式地提升各类视觉任务的性能，为实际工程系统提供即插即用的特征增强工具，具有广泛的应用前景。

\section{研究现状与挑战}

生成式模型旨在学习数据的底层分布$P(X)$或联合分布$P(X,Y)$，进而合成符合该分布的新样本。经典的生成式模型包括生成对抗网络（GAN）\cite{goodfellow2014generative,karras2019style,mirza2014conditional}和变分自编码器（VAE）\cite{kingma2013auto,van2017neural,zhao2017learning}等，它们通过学习隐空间表示捕获数据的生成过程，在图像合成、风格迁移等任务中取得了显著进展。近年来，扩散模型作为一类新兴的生成式框架引起了广泛关注。Sohl-Dickstein等人\cite{sohl2015deep}于2015年最早提出扩散概率模型的基本框架，通过迭代去除训练图像中的高斯噪声来建模数据分布。2020年，Ho等人\cite{ho2020denoising}提出去噪扩散概率模型（DDPM），通过精心设计的噪声调度策略和简化的训练目标，使扩散模型在图像生成质量上首次超越GAN，成为生成式建模的主流方法。此后，扩散模型在多种生成任务中持续展现出卓越能力：Stable Diffusion\cite{rombach2022high}通过在隐空间进行扩散过程实现了高分辨率图像合成，DALL·E\cite{ramesh2021zero}系列实现了零样本文本到图像生成，Imagen\cite{saharia2022photorealistic}通过结合深度语言理解实现了逼真的图像生成。此外，扩散模型还被拓展至视频生成\cite{ho2022imagen}和分子生成\cite{hoogeboom2022equivariant}等领域。在推理加速方面，DDIM\cite{song2020denoising}提出非马尔可夫扩散过程以减少采样步数，渐进式蒸馏\cite{salimans2022progressive}和DPM-Solver\cite{lu2022dpm}等方法进一步提升了采样效率。这些成果充分表明，扩散模型具备优良的数据分布建模能力和强大的噪声抑制能力\cite{10419041}。

尽管扩散模型在生成领域取得了巨大成功，但其在判别式任务中的探索仍然相对有限。与生成式模型通过建模联合分布$P(X,Y)$来合成样本不同，判别式模型关注的是条件分布$P(Y|X)$的估计，其中$Y$可以是分类标签、检测框或分割掩码等任务特定输出。近年来，一些研究开始尝试将扩散机制引入特定的判别式任务。DiffusionDet\cite{chen2023diffusiondet}将目标检测建模为从噪声框到目标框的去噪扩散过程，通过在检测框空间施加前向扩散并训练去噪网络逐步恢复精确的目标框坐标，在与成熟检测器的比较中取得了具有竞争力的表现。DiffSeg\cite{tian2023diffuse}利用预训练的稳定扩散模型实现了无监督零样本图像分割，通过迭代合并基于KL散度的注意力图生成有效的分割掩码，无需额外训练或文本监督即可获得良好的分割效果。此外，CLIP-Diffusion\cite{li2023clip}等工作探索了结合预训练视觉-语言模型与扩散机制的零样本学习方案。然而，上述方法均存在两方面的共同局限：其一，方法针对特定任务进行定制化设计，依赖于特定的数据结构（如DiffusionDet操作的是噪声边界框，DiffSeg操作的是噪声初始化的分割图），难以直接迁移至其他判别式任务，缺乏通用性；其二，递归式的去噪推理过程需要多步迭代采样，引入了显著的额外计算开销，限制了实际部署中的推理效率。

与此同时，特征噪声问题广泛存在于各类判别式视觉任务之中，是制约模型性能进一步提升的共性瓶颈。在图像分类任务中\cite{deng2009imagenet,krizhevsky2012imagenet}，尤其是细粒度分类场景\cite{wah2011caltech,parkhi2012cats,nilsback2008automated}下，类间视觉差异细微，特征中的噪声容易导致类间特征边界模糊，显著影响分类精度。在目标检测任务中\cite{ren2015faster,he2017mask,carion2020end}，背景杂乱、目标遮挡和尺度变化等因素引入的噪声会干扰区域特征的判别性，降低检测准确率。在图像分割任务中\cite{long2015fully,xie2021segformer,zhou2017scene}，像素级预测要求特征具有精细的空间对应关系，特征噪声容易导致分割边界不精确和小目标区域的错误预测。在图像检索任务中\cite{ye2021deep,bedagkar2014survey,lin2015deep}，以行人重识别为代表的细粒度检索需要跨摄像头匹配同一目标，面临姿态变化、光照差异、遮挡干扰等严峻挑战\cite{liao2015person,wang2020high,ni2023part}，对特征的身份判别性和鲁棒性要求极高。现有方法主要通过改进网络架构\cite{he2016deep,dosovitskiy2020image,liu2021swin}、设计注意力机制\cite{vaswani2017attention}或引入数据增强策略来缓解噪声影响，但尚缺乏一种从特征表示层面直接进行去噪增强的通用范式。

此外，知识蒸馏\cite{8395024}作为一种通过教师-学生架构进行模型压缩和知识传递的训练范式，近年来得到了广泛应用。在该框架中，参数量较大、表达能力较强的教师模型为轻量化的学生模型提供软监督信号，使学生模型在保持高效的同时获得接近教师模型的性能。在扩散模型的语境下，知识蒸馏同样发挥着重要作用：渐进式蒸馏\cite{salimans2022progressive}通过将多步采样过程压缩为更少的步数来加速扩散模型推理。这启发了一个重要思路——当去噪模块由于参数融合的需要被设计为轻量化结构时，可以借助更具表达能力的教师网络在训练阶段提供更精确的噪声分布估计，通过蒸馏将教师的去噪知识传递给轻量化的学生模块，从而在不增加推理开销的前提下提升去噪质量。上述各方面的研究现状将在第二章进行更为详尽的综述。

综合以上分析，当前将扩散模型的去噪能力应用于判别式表征学习面临以下三方面的关键挑战：

挑战一：缺乏通用的特征去噪框架。现有的扩散式判别方法均针对特定任务的特定数据结构设计去噪过程，无法通用地增强骨干网络提取的特征表示。亟需一种与任务无关的通用特征去噪范式，使生成式去噪能力能够无缝嵌入各类判别式框架之中。

挑战二：去噪过程引入不可忽视的推理开销。标准扩散模型的去噪过程采用递归式多步迭代，每一步都需要完整的网络前向计算，这在推理阶段引入了显著的额外计算开销。对于实时性要求严格的实际部署场景（如自动驾驶、视频监控），这种开销通常是不可接受的。如何实现去噪增强与推理零开销的兼顾，是一个核心技术难题。

挑战三：轻量化去噪模块的能力瓶颈。若为了实现推理零额外开销而将去噪模块设计为轻量化结构以便与骨干网络参数融合，则其有限的参数量会制约对复杂高维噪声分布的建模能力，导致去噪效果受限。如何在保持参数可融合性的同时提升去噪模块的有效去噪能力，需要新的训练策略和架构设计。

\section{研究内容与主要贡献}

针对上述挑战，本文以``基于扩散模型的图像判别任务特征增强''为核心研究主题，提出了两种互补的特征增强算法，系统性地探索了生成式去噪能力在判别式表征学习中的有效利用。

% 【建议插入图片：本文研究整体框架图】
% 内容：展示本文两个核心算法之间的关系和演进逻辑。
% 左侧：面向无额外推理成本的特征增强算法（第三章）——联合特征提取与去噪 + 参数融合
% 右侧：基于模型蒸馏的特征增强算法（第四章）——在左侧基础上增加教师-学生蒸馏和多步去噪融合
% 中间用箭头标注"增强与拓展"的关系，底部统一标注"推理零额外开销"。

研究内容一：面向无额外推理成本的特征增强算法（第三章）。针对挑战一和挑战二，本文提出了一种将扩散模型的去噪过程融入判别式骨干网络的通用特征增强方法。该方法的核心思想是：将骨干网络中级联的$N$个嵌入层（如卷积层或多头注意力层）重新解释为$T$步递归去噪层，使得特征提取与特征去噪在网络的逐层传播中同步完成。在此基础上，本文设计了一种参数融合机制，将训练好的去噪模块的参数直接合并至对应的嵌入层参数中，并从理论上证明了融合前后的数学等价性，从而在推理阶段完全消除了去噪过程带来的额外计算开销。该方法作为一种无标签增量式特征增强算法，可以即插即用地应用于各种预训练骨干网络和下游判别任务。

研究内容二：基于模型蒸馏的特征增强算法（第四章）。针对挑战三，本文在上述基础框架之上提出两个协同增强模块。第一，针对轻量化去噪模块建模能力不足的问题，引入教师-学生蒸馏框架：在训练阶段使用一个表达能力更强的教师网络来指导轻量化的学生去噪模块，教师网络提供更丰富、更具判别性的噪声分布估计，帮助学生模块更准确地捕捉结构化噪声模式，提升去噪方向的校准精度。第二，针对有限去噪深度下单步去噪存在的离散化误差问题，提出多步去噪融合策略：包括多轮单层去噪（将单个去噪层的多步迭代通过递推公式折叠为一次线性运算）和多步去噪蒸馏（借鉴DDIM\cite{song2020denoising}式跳步采样思想，将教师的多步去噪轨迹蒸馏至学生的单步去噪中）。上述增强模块仅在训练阶段引入额外计算，推理阶段仍然保持零额外开销的特性。

本文的主要贡献总结如下：

（1）提出了一种通用的基于扩散去噪的判别式特征增强框架，将骨干网络的$N$个级联嵌入层重新解释为$T$步递归去噪步骤，实现了特征提取与特征去噪的同步进行，生成更具判别性的特征表示。

（2）设计了特征提取与去噪参数融合策略，将去噪模块参数合并入嵌入层参数中，并从理论上证明了融合前后的数学等价性，实现了推理阶段零额外计算开销的特征增强。

（3）引入教师-学生蒸馏框架指导轻量化去噪模块的训练，通过更具表达能力的教师网络提供精确的噪声分布估计，在不增加推理开销的前提下显著提升了去噪模块的去噪能力。

（4）提出了多轮单层去噪与多步去噪融合两种增强策略，将多步去噪轨迹压缩为等效的单步操作，在有限去噪深度下增强了去噪性能和特征可分性。

（5）在行人重识别（Market-1501、DukeMTMC-reID、MSMT17、CUHK-03、VehicleID）、图像分类（ImageNet、CUB200、Oxford-Pet、Flowers）、目标检测（COCO）和图像分割（ADE20K）等多类视觉判别任务上进行了广泛实验，验证了所提方法在CNN架构（ResNet）和Transformer架构（ViT、Swin、VMamba）上的一致性能提升。

\section{论文结构安排}

本文围绕``基于扩散模型的图像判别任务特征增强算法研究''这一主题，从理论分析、算法设计到实验验证逐步展开，共分为五章。各章内容安排如下：

第一章为绪论。\\首先介绍图像判别式任务与表征学习的研究背景，阐述扩散模型在特征增强方面的应用意义；随后综述扩散模型、判别式表征学习和知识蒸馏等领域的研究现状，分析当前面临的三个关键挑战；在此基础上明确本文的两项研究内容并总结主要贡献；最后概述论文的整体组织结构。

第二章为相关工作。\\系统综述与本文研究密切相关的基础理论与技术方法。具体包括：生成式扩散模型的基本原理、前向扩散与反向去噪的数学框架及其发展历程；卷积神经网络、Vision Transformer等判别式表征学习的代表性骨干网络及其演进脉络；扩散模型在目标检测、图像分割等判别任务中的已有探索及其局限性；DDIM、渐进式蒸馏、DPM-Solver等高效扩散采样方法；以及教师-学生知识蒸馏的基本范式与典型应用。本章为后续两章的方法设计提供理论基础与技术参照。

第三章为面向无额外推理成本的特征增强算法研究。\\针对缺乏通用特征去噪框架和去噪推理开销两大挑战，本章提出将骨干网络中级联的嵌入层重新解释为递归去噪层的统一框架，实现特征提取与特征去噪的同步进行。在此基础上，设计参数融合机制将去噪模块参数合并入嵌入层参数中，并从理论上严格证明融合前后的数学等价性，从而在推理阶段完全消除额外计算开销。本章还详细介绍了实验设置与评价指标，并通过在行人重识别、图像分类、目标检测和图像分割等多类判别式任务上的对比实验，全面验证了所提方法的有效性、通用性和效率优势。

第四章为基于模型蒸馏的特征增强算法研究。\\针对第三章中轻量化去噪模块因参数量有限而导致建模能力受限的问题，本章首先分析从参数融合到蒸馏增强的问题演进逻辑，随后提出两个协同增强模块：一是引入教师-学生蒸馏框架，利用高容量教师网络在训练阶段指导轻量化学生去噪模块，提升噪声分布估计的精度；二是提出多步去噪融合策略，包括多轮单层去噪和多步去噪蒸馏两种具体方案，在有限去噪深度下增强去噪效果。本章通过详尽的消融实验验证各模块的独立贡献与协同效应，并通过跨任务、跨架构的泛化实验进一步确认方法的鲁棒性。

第五章为结论与展望。\\总结本文在基于扩散模型的判别式特征增强方面的主要工作与创新点，客观分析现有方法在适用场景、计算资源和理论深度等方面的局限性，并对跨模态特征去噪、自适应去噪调度以及与大规模预训练模型结合等未来研究方向进行展望。